\documentclass{article}
\usepackage{amsmath,amssymb,amsthm}
\usepackage{algorithm}
\usepackage[noend]{algpseudocode}
\usepackage{bm}
\usepackage{enumitem}
\usepackage{hyperref}
\usepackage{geometry}
\geometry{margin=1in}
\newtheorem{theorem}{Theorem}
\newtheorem{lemma}{Lemma}
\newtheorem{assumption}{Assumption}
\newtheorem{corollary}{Corollary}
\newcommand{\E}{\mathbb{E}}
\newcommand{\R}{\mathbb{R}}
\newcommand{\norm}[1]{\left\|#1\right\|}
\newcommand{\ip}[2]{\left\langle #1,#2\right\rangle}
\newcommand{\grad}{\nabla}
\newcommand{\clip}[1]{\operatorname{Clip}\!\left(#1\right)}
\begin{document}

\section*{Algorithm Name}
\textbf{Differentially Private Stochastic Gradient Descent with Gaussian Noise (DP‑SGD)}  

The algorithm performs standard SGD updates on a (convex, $L$‑smooth) loss while adding calibrated Gaussian noise to each mini‑batch gradient to guarantee $(\varepsilon,\delta)$‑differential privacy.

\section*{Mathematical Formulation}
Let $f:\R^{d}\to\R$ be the empirical risk
\[
f(w)=\frac{1}{n}\sum_{i=1}^{n}\ell(w;z_i),\qquad 
\ell(\cdot;z_i)\;\text{convex and $L$‑smooth}.
\]

At iteration $t$ a mini‑batch $\mathcal{B}_t\subset\{1,\dots,n\}$ of size $b$ is sampled uniformly
(without replacement).  
Define the (clipped) stochastic gradient
\[
g_t\;=\;\frac{1}{b}\sum_{i\in\mathcal{B}_t}\clip{\grad\ell(w_t;z_i)}\;+\;\xi_t,
\qquad 
\xi_t\sim\mathcal{N}(0,\sigma^{2}\mathbf{I}_d).
\]
Clipping is performed component‑wise to a radius $C>0$:
\[
\clip{\grad\ell(w;z)}\;=\;
\grad\ell(w;z)\cdot\min\!\Bigl\{1,\frac{C}{\norm{\grad\ell(w;z)}}\Bigr\}.
\]

The update rule is
\[
\boxed{w_{t+1}=w_t-\eta_t g_t},
\]
where $\eta_t>0$ is the learning‑rate schedule (e.g.\ $\eta_t=\eta$ or $\eta_t=\frac{\eta_0}{\sqrt{t}}$).

Privacy accounting: each iteration releases a Gaussian‑noised gradient with variance
$\sigma^{2}=\frac{2C^{2}\log(1.25/\delta)}{b^{2}\varepsilon^{2}}$
(standard Gaussian mechanism).  By the moments accountant (or advanced composition) the whole run of $T$ iterations satisfies $(\varepsilon,\delta)$‑DP.

\section*{Pseudocode}
\begin{algorithm}[H]
\caption{DP‑SGD}
\begin{algorithmic}[1]
\State \textbf{Input:} data $\{z_i\}_{i=1}^{n}$, loss $\ell$, clipping norm $C$, noise scale $\sigma$, batch size $b$, learning‑rate schedule $\{\eta_t\}_{t=0}^{T-1}$, iterations $T$.
\State Initialize $w_0\in\R^{d}$.
\For{$t=0$ \textbf{to} $T-1$}
    \State Sample a mini‑batch $\mathcal{B}_t$ of size $b$ uniformly at random.
    \State Compute per‑example gradients $g_i^{(t)}=\grad\ell(w_t;z_i)$ for $i\in\mathcal{B}_t$.
    \State Clip: $\tilde g_i^{(t)}=\clip{g_i^{(t)}}$.
    \State Average: $\bar g_t=\frac{1}{b}\sum_{i\in\mathcal{B}_t}\tilde g_i^{(t)}$.
    \State Sample noise $\xi_t\sim\mathcal{N}(0,\sigma^{2}\mathbf{I}_d)$.
    \State Form noisy gradient $g_t=\bar g_t+\xi_t$.
    \State Update $w_{t+1}=w_t-\eta_t g_t$.
\EndFor
\State \textbf{Output:} $w_T$ (or averaged iterate $\bar w_T=\frac{1}{T}\sum_{t=0}^{T-1}w_t$).
\end{algorithmic}
\end{algorithm}

\section*{Dry Run Example}
Consider the simple 1‑dimensional quadratic
\[
f(w)=(w-3)^2,\qquad\grad f(w)=2(w-3).
\]
No mini‑batching is needed; we set $C=\infty$ (no clipping) and $\sigma=0$ (pure SGD) to illustrate the deterministic part of the update.

\begin{center}
\begin{tabular}{c|c|c|c}
$t$ & $w_t$ & $g_t=\grad f(w_t)$ & $w_{t+1}=w_t-\eta g_t$ \\ \hline
0 & $0$ & $2(0-3)=-6$ & $0-0.1(-6)=0.6$ \\
1 & $0.6$ & $2(0.6-3)=-4.8$ & $0.6-0.1(-4.8)=1.08$ \\
2 & $1.08$ & $2(1.08-3)=-3.84$ & $1.08-0.1(-3.84)=1.464$ \\
\end{tabular}
\end{center}
Corresponding objective values: $f(0)=9$, $f(0.6)=5.76$, $f(1.08)=3.6864$, $f(1.464)=2.3507$. The iterates move toward the minimizer $w^{\star}=3$.

\newpage
\section*{Assumptions}
\begin{assumption}[Convexity]\label{ass:convex}
For all $w,w'\in\R^{d}$,
\[
f(w')\ge f(w)+\ip{\grad f(w)}{w'-w}.
\]
\end{assumption}
\begin{assumption}[L‑smoothness]\label{ass:smooth}
There exists $L>0$ such that for all $w,w'\in\R^{d}$
\[
f(w')\le f(w)+\ip{\grad f(w)}{w'-w}+\frac{L}{2}\norm{w'-w}^{2}.
\]
\end{assumption}
\begin{assumption}[Bounded (clipped) gradients]\label{ass:bddgrad}
For every example $z$ and every $w$, the clipped gradient satisfies
\[
\norm{\clip{\grad\ell(w;z)}}\le C.
\]
Consequently $\norm{\bar g_t}\le C$ for all $t$.
\end{assumption}
\begin{assumption}[Noise distribution]\label{ass:noise}
The added noise $\xi_t\sim\mathcal{N}(0,\sigma^{2}\mathbf{I}_d)$ is independent across iterations and independent of the sampled mini‑batch.
\end{assumption}
\begin{assumption}[Learning‑rate]\label{ass:lr}
The step sizes are non‑increasing and satisfy $\eta_t\le \frac{1}{L}$ for all $t$.
\end{assumption}

\section*{Main Convergence Theorem}
\begin{theorem}\label{thm:conv}
Let Assumptions~\ref{ass:convex}–\ref{ass:lr} hold. Define the averaged iterate
\[
\bar w_T:=\frac{1}{T}\sum_{t=0}^{T-1}w_t .
\]
If a constant step size $\eta_t=\eta\le\frac{1}{L}$ is used, then
\[
\E\!\bigl[f(\bar w_T)-f(w^{\star})\bigr]
\;\le\;
\frac{\norm{w_0-w^{\star}}^{2}}{2\eta T}
\;+\;
\frac{\eta C^{2}}{2}
\;+\;
\frac{\eta d\sigma^{2}}{2},
\tag{1}
\]
where the expectation is over the randomness of mini‑batch sampling and the Gaussian noise.
\end{theorem}

\section*{Theoretical Properties}
\begin{itemize}[leftmargin=*]
\item \textbf{Stability.} The bound (1) is decreasing in $T$ provided $\eta$ is chosen as $\eta=O(1/\sqrt{T})$, yielding the optimal $O(1/\sqrt{T})$ rate for convex problems.
\item \textbf{Privacy–utility trade‑off.} The noise variance $\sigma^{2}$ is dictated by the privacy budget $(\varepsilon,\delta)$. Larger $\sigma^{2}$ inflates the third term in (1), degrading utility.
\item \textbf{Comparison to vanilla SGD.} Setting $\sigma=0$ recovers the classic SGD bound 
$\frac{\norm{w_0-w^{\star}}^{2}}{2\eta T}+ \frac{\eta C^{2}}{2}$.
\item \textbf{Hyper‑parameter sensitivity.} The bound is linear in $C$ and $\sqrt{\sigma^{2}}$, suggesting that tighter clipping and smaller noise (i.e.\ larger $\varepsilon$) improve convergence.
\end{itemize}

\section*{Discussion}
The theorem shows that DP‑SGD preserves the $O(1/\sqrt{T})$ convergence rate of standard SGD up to an additive penalty proportional to the noise variance. By calibrating $\sigma$ according to the Gaussian mechanism, the algorithm satisfies $(\varepsilon,\delta)$‑DP for any prescribed privacy budget. The bound is non‑asymptotic and holds for any convex, $L$‑smooth loss, making it directly applicable to empirical risk minimisation in machine learning.

\newpage
\section*{Preliminaries (Definitions and Lemmas)}
\begin{lemma}[Three‑point identity]\label{lem:three}
For any $a,b\in\R^{d}$ and any $w\in\R^{d}$,
\[
\norm{a-w}^{2}= \norm{b-w}^{2}+2\ip{a-b}{b-w}+\norm{a-b}^{2}.
\]
\end{lemma}
\begin{proof}
Expand both sides using $\norm{x}^{2}= \ip{x}{x}$ and collect terms. \qedhere
\end{proof}

\begin{lemma}[Unbiasedness of clipped gradient]\label{lem:unbiased}
Conditioned on $w_t$, the expectation of the noisy gradient satisfies
\[
\E[g_t\mid w_t]=\E[\bar g_t\mid w_t]=\frac{1}{b}\sum_{i\in\mathcal{B}_t}\E\bigl[\clip{\grad\ell(w_t;z_i)}\mid w_t\bigr].
\]
\end{lemma}
\begin{proof}
Noise $\xi_t$ has zero mean and is independent of the batch, thus $\E[\xi_t]=0$. The remaining term is deterministic given the sampled batch, yielding the claim. \qedhere
\end{proof}

\section*{Main Proof (Step‑by‑Step Derivation)}
We prove Theorem~\ref{thm:conv}. Throughout the proof $w^{\star}$ denotes any minimiser of $f$.

\noindent\textbf{Step 1 (Distance recursion).}  
Using the update $w_{t+1}=w_t-\eta g_t$ and Lemma~\ref{lem:three} with $a=w_{t+1}$, $b=w_t$, we obtain
\[
\norm{w_{t+1}-w^{\star}}^{2}
 = \norm{w_t-w^{\star}}^{2}
   -2\eta\ip{g_t}{w_t-w^{\star}}
   +\eta^{2}\norm{g_t}^{2}.
\tag{2}\label{eq:rec}
\]

\noindent\textbf{Step 2 (Take conditional expectation).}  
Condition on $w_t$ and apply Lemma~\ref{lem:unbiased}:
\[
\E\!\bigl[\norm{w_{t+1}-w^{\star}}^{2}\mid w_t\bigr]
 = \norm{w_t-w^{\star}}^{2}
   -2\eta\ip{\E[g_t\mid w_t]}{w_t-w^{\star}}
   +\eta^{2}\E\!\bigl[\norm{g_t}^{2}\mid w_t\bigr].
\tag{3}
\]

\noindent\textbf{Step 3 (Decompose the inner product).}  
Because $f$ is convex (Assumption~\ref{ass:convex}),
\[
\ip{\grad f(w_t)}{w_t-w^{\star}}\;\ge\; f(w_t)-f(w^{\star}).
\tag{4}
\]
The clipped gradient is bounded by $C$ (Assumption~\ref{ass:bddgrad}), hence
\[
\bigl\|\clip{\grad\ell(w_t;z)}\bigr\|\le C\quad\Longrightarrow\quad
\|\bar g_t\|\le C.
\tag{5}
\]

\noindent\textbf{Step 4 (Bound the expected squared norm).}  
Write $g_t=\bar g_t+\xi_t$. Using independence of $\xi_t$ and the bound (5):
\[
\begin{aligned}
\E\!\bigl[\norm{g_t}^{2}\mid w_t\bigr]
 &= \E\!\bigl[\norm{\bar g_t+\xi_t}^{2}\mid w_t\bigr]  \\
 &= \norm{\bar g_t}^{2}+ \E\!\bigl[\norm{\xi_t}^{2}\bigr]    \qquad\text{(cross term zero)}\\
 &\le C^{2}+ d\sigma^{2}.
\end{aligned}
\tag{6}
\]

\noindent\textbf{Step 5 (Replace $\E[g_t\mid w_t]$ with true gradient).}  
Since clipping is a non‑expansive projection, we have for any $z$
\[
\ip{\clip{\grad\ell(w_t;z)}}{w_t-w^{\star}}
 \;\le\; \ip{\grad\ell(w_t;z)}{w_t-w^{\star}} .
\]
Averaging over the batch and taking expectation yields
\[
\ip{\E[g_t\mid w_t]}{w_t-w^{\star}}
 \;\le\; \ip{\grad f(w_t)}{w_t-w^{\star}} .
\tag{7}
\]

\noindent\textbf{Step 6 (Combine (3)–(7)).}  
Insert (4)–(7) into (3):
\[
\begin{aligned}
\E\!\bigl[\norm{w_{t+1}-w^{\star}}^{2}\mid w_t\bigr]
 &\le \norm{w_t-w^{\star}}^{2}
      -2\eta\bigl[f(w_t)-f(w^{\star})\bigr]
      +\eta^{2}\bigl(C^{2}+d\sigma^{2}\bigr).
\end{aligned}
\tag{8}
\]

\noindent\textbf{Step 7 (Take total expectation).}  
Removing the conditioning,
\[
\E\!\bigl[\norm{w_{t+1}-w^{\star}}^{2}\bigr]
 \le \E\!\bigl[\norm{w_t-w^{\star}}^{2}\bigr]
      -2\eta\E\!\bigl[f(w_t)-f(w^{\star})\bigr]
      +\eta^{2}\bigl(C^{2}+d\sigma^{2}\bigr).
\tag{9}
\]

\noindent\textbf{Step 8 (Summation over $t=0,\dots,T-1$).}  
Telescoping the left‑hand side yields
\[
\begin{aligned}
\E\!\bigl[\norm{w_T-w^{\star}}^{2}\bigr]
 &\le \norm{w_0-w^{\star}}^{2}
      -2\eta\sum_{t=0}^{T-1}\E\!\bigl[f(w_t)-f(w^{\star})\bigr]
      +T\eta^{2}\bigl(C^{2}+d\sigma^{2}\bigr).
\end{aligned}
\tag{10}
\]

\noindent\textbf{Step 9 (Rearrange).}  
Since the squared norm is non‑negative,
\[
2\eta\sum_{t=0}^{T-1}\E\!\bigl[f(w_t)-f(w^{\star})\bigr]
 \le \norm{w_0-w^{\star}}^{2}+T\eta^{2}\bigl(C^{2}+d\sigma^{2}\bigr).
\tag{11}
\]

\noindent\textbf{Step 10 (Average iterate).}  
Divide (11) by $2\eta T$ and use convexity of $f$:
\[
\begin{aligned}
\E\!\bigl[f(\bar w_T)-f(w^{\star})\bigr]
 &\le \frac{1}{T}\sum_{t=0}^{T-1}\E\!\bigl[f(w_t)-f(w^{\star})\bigr] \\
 &\le \frac{\norm{w_0-w^{\star}}^{2}}{2\eta T}
      +\frac{\eta}{2}\bigl(C^{2}+d\sigma^{2}\bigr).
\end{aligned}
\tag{12}
\]

Equation~(12) is exactly the bound (1) of Theorem~\ref{thm:conv}. \hfill\qedsymbol

\section*{Rate Derivation (With Constants)}
Choosing $\eta = \frac{R}{\sqrt{T(C^{2}+d\sigma^{2})}}$ where $R:=\norm{w_0-w^{\star}}$, we obtain
\[
\E\!\bigl[f(\bar w_T)-f(w^{\star})\bigr]
 \le \frac{R\sqrt{C^{2}+d\sigma^{2}}}{\sqrt{T}}.
\]
Thus the algorithm attains the optimal $O(1/\sqrt{T})$ rate for convex objectives, with a constant that grows linearly with the clipping radius $C$ and with $\sqrt{d}\,\sigma$ (the privacy‑induced noise level).

\section*{Special Cases}
\begin{enumerate}[leftmargin=*]
\item \textbf{No privacy noise ($\sigma=0$).} The bound reduces to the classical SGD guarantee:
\[
\E[f(\bar w_T)-f(w^{\star})]\le\frac{\norm{w_0-w^{\star}}^{2}}{2\eta T}+\frac{\eta C^{2}}{2}.
\]
\item \textbf{Strongly convex $f$ (parameter $\mu>0$).} Repeating the above steps with the stronger inequality
$\ip{\grad f(w_t)}{w_t-w^{\star}}\ge \mu \norm{w_t-w^{\star}}^{2}+f(w_t)-f(w^{\star})$
yields a linear convergence term plus an additive $\frac{\eta d\sigma^{2}}{2\mu}$ steady‑state error.
\item \textbf{Full‑batch DP‑GD (batch size $b=n$).} The clipping bound becomes $C\le \max_i\norm{\grad\ell(w;z_i)}$, and the noise variance scales as $\sigma^{2}\propto \frac{C^{2}}{n^{2}\varepsilon^{2}}$, improving the utility term by a factor $n$.
\end{enumerate}

\end{document}